\begin{empiricalfinding}[label={find:census}]{The full permutation census}
Every permutation containing an odd cycle lacks a fixed byte
(Supplement~2, Theorem~1); under the census protocol, all $29{,}295$ such
operators were classified genotype-alive. The theorem excludes a common fixed
byte for all elementary writes; the observed $100\%$ alive rate is an additional
dynamical result. Every all-even cycle type splits between live and dead
operators. Among those types, the alive fraction falls as the number of cycles
$k$ increases: $[8]$ ($k{=}1$) is $83\%$ alive, $[2\,6]$ and $[4\,4]$ ($k{=}2$)
$70$/$67\%$, $[2\,2\,4]$ ($k{=}3$) $49\%$, and $[2\,2\,2\,2]$ ($k{=}4$) $26\%$.
Overall, $92.5\%$ have a live genotype. The regime split is $82.7\%$ live/live,
$9.8\%$ live/dead, $6.0\%$ dead/dead, and $1.4\%$ dead/live, so all four regimes
occur among bijective operators.
\end{empiricalfinding}

\begin{empiricalfinding}[label={find:survey}]{The rotation survey}
Within the wrap-rotation subfamily, every non-trivial rotation is all-even and
so possesses fixed rules (Supplement~2, Theorem~1), yet under the
neighbour-blend coupling the tested rotations split by cycle structure
(Supplement~3, Figure~1 and Figure~\ref{fig:lambda}):
\begin{itemize}\setlength{\itemsep}{1pt}
\item the single 8-cycles---$\gcd(\text{offset},8)=1$, offset $\pm 1$ and
$\pm 3$---\emph{sustain} a high-diversity field, holding $\sim$30--40 distinct
rules through the longest runs we made ($400$ generations);
\item offset $\pm 2$ ($\gcd 2$, two $4$-cycles) and offset $\pm 4$ ($\gcd 4$,
four $2$-cycles) \emph{collapse} to one or two rules within a few dozen
generations.
\end{itemize}
The distinct-rule population curve (Figure~\ref{fig:lambda}) separates the two by
more than an order of magnitude in a single scalar. Possessing a fixed rule and
having the coupled field stay near it are different properties. The present
claim is limited to two seeds per rotation and treats population size as a
descriptive observable, not as an order parameter in the statistical-mechanics
sense.
\end{empiricalfinding}

\begin{empiricalfinding}[label={find:braid}]{A braid of two collapsing operators can sustain the field}
At width $80$, horizon $120$, and seed $0$, offset $+2$ and offset $+4$ each
have genotype change rate $0$ when held alone, while their alternating braid
has change rate $0.76$ (Figure~\ref{fig:braid}, top). The matched schedule
control offset $+2$/offset $-2$ remains dead (Figure~\ref{fig:braid}, bottom).
This establishes that two collapsing components can construct a sustaining
schedule. It does not establish the previously suggested block-system
criterion: all three even offsets preserve the evens/odds partition, and both
pairs generate the same translation subgroup $\langle {+2}\rangle$. The cause
of the contrasting braid outcomes remains to be identified.
\end{empiricalfinding}

\begin{empiricalfinding}[label={find:knob}]{The mixing fraction is a reproducible diversity control}
Let $\sigma$ be displayed-position offset $+1$. We ran periodic 12-step
schedules mixing the sustaining 8-cycle $\sigma^1$ with the collapsing
two-4-cycle operator $\sigma^2$, at width $80$ for $120$ generations and five
seeds per mixture. Mean late-window diversity rises from $1.0$ rule when the
$\sigma^1$ fraction is zero, through $16.8$ at fraction $0.25$, to $36.9$ at
fraction $5/12$; above that it remains in the range $34$--$41$ rather than
increasing monotonically (Figure~\ref{fig:knob}). The schedule fraction is thus
a tunable control for entry into a sustained high-diversity regime, with
saturation rather than a continuously linear response.
\end{empiricalfinding}

\begin{empiricalfinding}[label={obs:couplings}]{The rotation split occurs with and without neighbour blending}
Figure~\ref{fig:parity} uses stochastic elementary writes without neighbour
blending; Supplement~3, Figure~1 and Figure~\ref{fig:lambda} include the blend. Over the
displayed horizons, offset $+1$ remains active and offset $+2$ settles under
both protocols, although their diversity levels and transients differ. These
runs therefore do not identify neighbour blending as the cause of the split;
they show that fixed-byte existence alone does not determine finite-horizon
settling under either protocol.
\end{empiricalfinding}

\begin{empiricalfinding}[label={obs:metrics}]{The tested per-trajectory metrics do not separate the parity classes}
The classification is a conjugacy invariant of $\sig$---its cycle type---while
the complexity and compression measures we tried read a single rule's bits or one
space--time trajectory without explicitly encoding $\sig$. In our tests these
measures did not recover the parity classification: a population count reads a collapsed-but-maximally-selected
field as nearly dead; a frozen periodic genotype scores a lower statistical
complexity $C_\mu$ \parencite{crutchfield1989inferring} than Rule~110; gzip rates
a self-similar rule as noise. This is a negative result about the tested
summaries, not an impossibility theorem for trajectory-based inference.
\end{empiricalfinding}

\begin{empiricalfinding}[label={find:crossover}]{No finite-size evidence of a critical point over the tested range}
The natural control parameter is the \emph{propagator fraction} $q$ ($q = 0$ all
blend, $q = 1$ all propagator, offset $+2$). If the sustained phase sat at a
critical point in $q$, a finite-size scan would show three things as the lattice
grows: the transition in genotype activity $a_G(q)$ \emph{sharpening}, its width
shrinking toward zero; a susceptibility proxy peaking at a \emph{fixed} interior
value of $q$, the peak growing without bound; and Binder cumulants for different
$L$ crossing at that same value. Across widths $L = 30$--$240$ (32 seeds each)
we observe none of the three.

First, $a_G(q)$ rises \emph{smoothly}, well described by a logistic crossover
($R^2 > 0.99$, Figure~\ref{fig:phase}a), and its width does not shrink over the
tested sizes ($w \approx 0.13$, $w \sim L^{-0.04}$)---the transition is not
sharpening. Second, the size-scaled run-to-run spread $L\,\mathrm{Var}(a_G)$
(Figure~\ref{fig:phase}b), used as a finite-size susceptibility proxy, does
\emph{peak} inside the crossover band and does grow (${\sim}5.7\times$ from the
smallest width to the largest), but its location drifts rather than converging
($q=0.75,0.50,1.00,1.00$ as $L$ increases), and for the two largest sizes the
maximum lies at the sampled boundary $q=1$, so no interior peak is resolved at
all. Third, there is no Binder-cumulant crossing. The scan therefore provides no
finite-size evidence for a sharpening transition over $L=30$--$240$. These results support a broad crossover, plausibly
influenced by seed-dependent residence times, rather than a confirmed critical
point; they do not exclude different scaling outside the tested range.
\end{empiricalfinding}

\begin{empiricalfinding}[label={find:river}]{The river sustains structured phenotype}
In all six tested seeds at $L=80$, $T=120$, the river produces coherent bands
of alternating phenotype that drift through a disordered background. In the
larger pinned run at $L=240$, $T=360$ (Figure~\ref{fig:river}), those structures
persist while branching, merging, and reorganising across the full spacetime
field. This is a bounded six-seed observation plus one larger representative
run, not a claim that the structure is an attractor or survives arbitrarily
long.
\end{empiricalfinding}

\begin{empiricalfinding}[label={find:causal}]{Phenotype-to-genotype feedback roughly doubles causal reach}
All runs use $L = 80$, $T = 120$, $t^{*} = 60$, with damage measured on the
\emph{phenotype} layer as the number of cells differing between branches and
distance as the shorter circular displacement. At $dt \le 10$ the two branches
are empirically indistinguishable. By $dt = 59$, in a single-seed sweep over all
$80$ sites, the live construction has damaged $12.97$ cells against the
control's $5.51$ and spread $16.95$ against $7.46$
(Figure~\ref{fig:river-perturbation} plots that sweep). Replication uses a
separate design---six seeds, $20$ evenly spaced sites each, aggregated to one
mean per seed so that the \emph{seed} is the unit of replication. The live
construction exceeds its control in $6/6$ seeds, per-seed differences
$+1.90$, $+7.95$, $+6.65$, $+6.60$, $+6.05$, $+0.85$ cells, mean $+5.00$
(sd $2.90$). Because $n = 6$ makes a $t$-statistic fragile we quote the exact
sign-flip permutation test over all $2^{6}$ assignments: the observed mean is
the extremum, one-sided $p = 1/64 \approx 0.016$ (paired $t(5) = 4.23$ for
comparison). Perturbations do not spread uniformly: in the single-seed sweep
$22$ of $80$ sites are dead ends, and they form a contiguous region.

One qualification on what the control cuts. \texttt{run-river-ablated-from}
reads a phenotype frozen at the start of each run segment, so it cuts the
\emph{dynamic} edge but still presents the genotype with a static field that,
in the perturbed branch, already contains the flip. The $5.51$ is therefore
feedback-frozen, not feedback-absent; with the frozen reference supplied
externally and identical in both branches, reach falls further, to $1.28$
(\findref{find:gain}). The doubling quoted here is the effect of making the
read \emph{current}, which is the smaller half of the total coupling effect.
\end{empiricalfinding}

\begin{empiricalfinding}[label={find:ladder}]{Causal reach on a scale calibrated against elementary rules}
On a third design---four seeds, ten sites each, damage recorded at every
$dt$---width grows as $dt^{0.64}$ for the live construction against $dt^{0.43}$
for its control, with centroid displacement scaling as $dt^{0.56}$ against
$dt^{0.43}$. Exponents are least-squares fits over $dt \in [5,59]$ on
site-averaged curves, conditional on damage existing---a run whose perturbation
has died contributes no width and is excluded rather than entered as zero---and
we quote no interval on them
(Figure~\ref{fig:river-centroid}); these are reported descriptively and we draw
no inference from them about gliders or particles, for a reason given below.
Placed on one scale against elementary rules run under an identical protocol
(mean damaged cells at $dt=59$): rule $0$ $0.00$; rule $204$ $1.00$; rule $90$
$8.00$; \textbf{the construction $12.97$}; rule $54$ $15.05$; rule $110$
$20.23$; rule $30$ $37.50$. The instrument returns exactly zero where nothing
propagates, and the construction propagates at a magnitude comparable to the
elementary rules that do.
\end{empiricalfinding}

\begin{empiricalfinding}[label={find:architecture}]{The order parameter is the strength of the phenotype-to-genotype coupling}
Run under the protocol above---$L=80$, $t^{*}=60$, one phenotype flip, damaged
phenotype cells at $dt=59$---the constructions separate by \emph{coupling}, not
by which parameter they vary. Where the genotype update never reads the
phenotype, nothing clears the bottom of the complex band: the bare operators
($P_a$ $0.03$, rot$+1$ $2.05$), braided writings ($0.12$, $5.20$), spatial
niches ($0.50$, $0.72$), asynchronous refuges ($0.47$, $0.75$), replacement
mutation ($0.23$, $0.85$), the preserving limit ($1.10$), and blend strength
across its whole range---$1.00$ giving $0.03$, $0.70$ giving $0.55$, $0.35$
giving $3.38$, and $0.00$, the most propagating feedforward configuration we
found, giving $8.15$, level with rule $90$. Most sit below rule $204$ ($1.00$),
the frozen identity.

The coupling is not merely present or absent; it is a dial, and reach is graded
and monotone along it. Conservative transport whose swap probability is gated on
the local phenotype interface rises from $1.10$ at rate $0$ through $4.25$ at
$0.10$, $9.43$ at $0.20$, $14.45$ at $0.35$, $19.27$ at $0.50$ and $26.05$ at
$0.75$, crossing rule $90$ between rates $0.10$ and $0.20$ and passing rules
$110$ and $54$ before saturating. The river construction sits at $12.97$ on the
same scale, and cutting its single edge returns it to $5.51$, back among the
feedforward constructions.

That this is the \emph{coupling} and not merely rule mobility is established by
an ungated control. Replacing the phenotype-gated swap probability with a
constant one leaves the mean number of swaps unchanged---for a balanced field
the gated mean is exactly the rate---and leaves the dynamics bijective and the
rule histogram invariant, so the two differ only in whether the genotype can see
the phenotype. The ungated control never leaves the ordered band: $1.65$, $1.68$,
$1.77$, $2.77$ and $4.05$ at rates $0.20$ through $1.00$, against $9.43$ to
$25.38$ for the gated version, a factor of $5.7$ to $10.9$ at matched mobility
(Figure~\ref{fig:regime-placement}).

An order parameter therefore exists, and the classical picture is recovered in
it: reach is a graded, monotone function of coupling strength with a locatable
crossing into the complex band. What is unusual is only \emph{which} coordinate
it is. Every parameter internal to the update---$\lambda$, propagator fraction,
sustained diversity, blend strength, mutation rate, refuge probability, niche
width, and rule mobility itself---is silent or nearly so. The coordinate that
governs is the gain of the loop from the phenotype back to the genotype, which
conventional scans do not sweep because architecture is usually treated as a
property a model has or lacks rather than as a quantity.
\end{empiricalfinding}

\begin{empiricalfinding}[label={find:gain}]{Coupling is a dial, not a switch, and reach is monotone along it}
Both constructions that read the phenotype admit a gain parameter, and in both
causal reach is graded and monotone in it under the protocol of
\findref{find:ladder}.

Figure~\ref{fig:gain-curves} plots both dials.
For conservative transport the gain is the phenotype-gated swap rate. Reach
rises $1.10$, $2.17$, $4.25$, $9.43$, $14.45$, $19.27$, $26.05$, $25.38$ at
rates $0$, $0.05$, $0.10$, $0.20$, $0.35$, $0.50$, $0.75$, $1.00$: it crosses
rule $90$ between rates $0.10$ and $0.20$, passes rules $110$ and $54$, and
saturates past $0.75$.

For the river the gain is the \emph{currency} of the four context bits the
genotype step reads. Each cell at each step reads the live phenotype with
probability $\gamma$ and a frozen reference otherwise, so $\gamma$ is the
fraction of the genotype's view of the phenotype that is causally current. Over
three seeds and eighty sites, reach rises $1.28$, $3.00$, $3.85$, $3.43$,
$4.63$, $5.00$, $5.77$, $7.38$, $12.39$ across $\gamma = 0$ to $1$ in eighths,
monotone within one standard error, with standard errors $0.19$ to $0.57$.

Three properties make the river dial an instrument rather than a knob. It is
anchored at both ends: $\gamma = 1$ reproduces the live river at $12.97$
exactly. The random tape is preserved---the propagator draw is consumed once per
cell at every $\gamma$, with gate coins from a separate stream re-seeded
identically in both branches, so the gate injects no divergence of its own. And
because the frozen field is a real phenotype, marginals and spatial structure
are matched at every $\gamma$; only causal currency varies.

The two curves have opposite curvature. Transport is concave, buying most of its
reach in the first fifth of the dial. The river is convex: $45$\% of its span
arrives in the last eighth, and only $\gamma = 1$ clears rule $90$, so one stale
read in eight costs $40$\% of the reach. A plausible reading is that the river's
first-match fires on a four-bit conjunction, so a stale context does not degrade
the match but breaks it; we have not tested this, and gating the four bits
independently rather than together would separate it.

A one-seed pilot showed a non-monotonic dip near $\gamma = 0.5$ that did not
survive three seeds. We record it because we nearly believed it.
\end{empiricalfinding}

\begin{empiricalfinding}[label={find:diversity-axis}]{Sustained diversity does not determine causal reach; conservative transport breaks the apparent interior optimum}
The original $76$ configurations produce an apparent rise and fall: binned
medians rise from $0.10$ cells below $40$ sustained rules to $2.34$ at
$125$--$135$, then fall to zero above $135$. That descent is completely
confounded. Every original configuration above $135$ obtains its diversity
either by replacement mutation or by holding the genotype fixed. The observed
failures are real---mutation erases the lesion and holding prevents its
spread---but they do not show that high diversity suppresses propagation.

We therefore add a conservative eighth mechanism. Each run begins from a
balanced empirical rule measure with support
$32,64,96,128,160,192,224,$ or $256$. Alternating layers of disjoint local
pairs exchange their rules with a probability controlled by a transport rate
and the local phenotype interface. These exchanges are bijections: no rule is
created or destroyed, and the complete genotype histogram and its Shannon
diversity remain invariant at every checked step. A local transposition supplies
the perturbation while preserving that same measure in both branches. We run
five transport rates, six seeds, and eight perturbation sites for $40$ response
steps, giving $1{,}920$ paired responses.

Figure~\ref{fig:diversity-axis} reports causal amplification---final genotype
damage divided by the two initially exchanged sites---so the conservative
responses share a scale with the original single-site survey. With transport
zero, amplification is exactly $1.00$ at every diversity. At transport rates
$0.25$, $0.50$, $0.75$, and $1.00$, the ranges across the full diversity axis
are respectively $6.63$--$8.26$, $14.11$--$15.31$, $18.75$--$21.26$, and
$20.78$--$22.08$. There is no descending limb: within each transport rate the
dependence on diversity is approximately flat or slightly positive, while
changing transport changes reach by more than an order of magnitude.

Sustained rule diversity is therefore not an order parameter for causal reach.
Mutation and holding are distinct suppression mechanisms, but they are not
flanks of a common diversity optimum. Structured genotype transport is the
control exposed by this experiment.
\end{empiricalfinding}
